\documentclass[a4paper,12pt]{article}
\usepackage{ngerman}
\usepackage{amsfonts}
\usepackage{amsmath}
\author{Alexander Engelhardt}
\date{\today}

\begin{document}
	\paragraph{Konvergenzenkalkulationen für $\lim_{n\to\infty}\{f(n) \mid n \in \mathbb{N} \setminus{\{0\}}\}$:}
	\begin{enumerate}
		\item 
		\[
			O(n^{2.5}) \rightarrow \lim_{n\to\infty}n^{2.5} = \infty
		\]\\
		$f'(n) = 2.5n^{1.5}$
		\item 
		\[
			O(n^2\log n) \; \longrightarrow\; \lim_{n\to\infty}
			\overbrace{%
				n^2 \,
				\overbrace{\log n}^{\text{\small$\infty$}}
				}^{\text{\small$\infty$}} 
				\; \Leftrightarrow\; \infty\ 
		\]\\
		$f'(n) = 2n\log n + \dfrac{1}{n}n^2$
		\item 
		\[
			O(\ln{(\ln{n})}) \; \longrightarrow \lim_{n\to\infty}\ln{(\ln{n})} \leftarrow O(\ln n) = O(\log n) \leftrightarrow \ln \infty = \infty 
		\]\\
		$f'(n) = \dfrac{1}{\ln n}\cdot\dfrac{1}{n}$
		\item $O(\log{n}) \rightarrow \lim_{n\to\infty}\log{n} = \infty$, s.o.\\
			  $f'(n) = f'(\ln n) = \dfrac{1}{n}$
		\item $O((\log{n})^3 \rightarrow \lim_{n\to\infty}(\log{n})^3 = \infty$, ebd.\\
			  $f'(n) = (\dfrac{1}{n})^3$
		\item $O(7) \rightarrow \lim_{n\to\infty}7 = 7 \Rightarrow$ konstante Komplexit"at ohne Wachstum\\
		$f'(n) = 0$
		\item $O(\sqrt{n}) \rightarrow \lim_{n\to\infty}\sqrt{n} = \infty$\\
			  $\sqrt{n} \leftrightarrow n^{\dfrac{1}{2}} \Rightarrow f'(n) = \dfrac{1}{2\sqrt{n}}$
		\item 
		\[
			O(n^2 + (\log{n})^3) \; \longrightarrow \lim_{n\to\infty}
			\overbrace{%
				n^2 \,
				+ \,
				\overbrace{(\log{n})^3)}^{\text{\small$\infty$}}
			}^{\text{\small$\infty$}} \; \leftrightarrow \infty
		\]\\
		$f'(n) = 2n + (\dfrac{1}{n})^3$
		\item 
		\[
			O(n!) \rightarrow \lim_{n\to\infty}n! = \prod_{n=1}^{k}n \leftrightarrow n \cdot (n-1) \cdot (n-2) \cdot ... \cdot 2 \cdot 1 \Rightarrow n! = n(n-1)! = \infty
		\]\\
		\text{$n!$ ist an sich nicht (sinnvoll) ableitbar, aber man kann folgende Annahme treffen:}
			\textbf{$n!$ wächst asymptotisch deutlich schneller als alle anderen hier genannten Funktionen, außer $n^n$ (s.u.)}
		\item $O(n^n) \rightarrow \lim_{n\to\infty}n^n = \infty$\\
			  $f'(n) = n^n(\ln n + 1) \Leftrightarrow n^n \cdot \ln n + 1 \cdot n^n$\\
			  $\Rightarrow \forall f(n) = c + (a\log_b n)^k :  a,b,c,k > 0 \land c \ne n^n \Rightarrow f(n) \in o(n!)$
		\item $O(2^n) \rightarrow \lim_{n\to\infty}2^n = \infty$\\
			  $f'(n) = \ln{(2)} \cdot 2^n$
		\item $O(\sqrt[5]{n}) \rightarrow \lim_{n\to\infty}\sqrt[5]{n} = \infty$\\
			  $f'(n) = \dfrac{1}{5\sqrt[5]{n}} \Rightarrow \sqrt[5]{n} \in o(\sqrt{n})$
		\item $O(176n^2) \rightarrow \lim_{n\to\infty}176n^2 = \infty$\\
			  $f'(n) = (176 \cdot 2)n = 352n$
	\end{enumerate}
	\paragraph{Komplexitätsklassen asymptotisch geordnet von langsam nach schnell wachsend (anhand der Steigung $f'(n)$):\\}
	\textbf{n.b: $\exists c \in \mathbb{N} > 0 \Rightarrow O(1), 1 \subset o(\log n) \subset o(n) \subset o(n \log n) \subset o(n^2) \\\subset o(n^3) \subset ...  \subset o(n^k) \subseteq o(2^n) \subseteq o(n!) \subseteq o(n^n)$}
	\begin{enumerate}
		\item $O(7) \Rightarrow f'(7) = 0$
		\item $O(\ln{(\ln n)}) \Rightarrow \ln{(\ln n)} \subseteq o(\ln n)$
		\item $O(\log n)$
	\end{enumerate}
\end{document}